\scanpage{103}
\hrulefill

Let $X/K$, and let $T$ be a subset of
\inserted{\tikz[baseline=0pt,x=1em,y=1ex]{
  \node[inner sep=0pt] at (0,1.8) {$X=$};
  \draw[line width=.4pt] (-.65,1.0) .. controls (-.3,.45) and (.45,.55) .. (.45,-.1) -- (.62,-1.8);
}}$X_{\mathbb C}(\mathbb C)$, local\textsuperscript{ly} closed 1-connected. Let $X'$ \deleted{be an algebraic variety}\inserted{$-X$ be a} $K$-morphism, which over \deleted{the part of} a neighborhood of $T$ is finite étale. Then $X'\vert T=X'_T$ decomposes into a finite number of copies of $T$ -- there is an equivalence relation in $X'_T$, which is called the $T$-adjacency of points of $X'$ over $T$. We are given \deleted{$I\subset X'_T$} $x',y'\in X'_T$
% The following small baseline trace is distinct from the overwork on the inserted T.
\unsure{\tikz[baseline=0pt,x=1ex,y=1ex]{
  \draw[line width=.45pt] (.6,.8) -- (.32,.58)
    .. controls (-.05,.55) and (-.18,.2) .. (.05,.08)
    .. controls (.38,-.02) and (.4,.3) .. (.28,.44)
    .. controls (.35,-.1) and (.15,-.5) .. (.05,-.7);
}}
\inserted{$T$ adjacent;}
\margin{\unsure{assum\textsuperscript{ing}|assuming} $x',y'\in X'_{\overline K}$}
if $u\in\Gamma$, one cannot assert that \deleted{\unsure{they are}}\inserted{$u(x'),u(y')$} are still $T$ adjacent.
But if they are, one can look at $l^T_{y',x'}:x'\to y'$ \inserted{``$T$-horizontal''} in \deleted{$\unsure{\in}\,\unsure{\widehat{\Pi}_1(X'_{\overline K},x',y')}$} $\operatorname{Hom}_{\widehat{\Pi}_1(X'_{\overline K})}(x',y')$ and its transform, \deleted{\unsure{the} \illegible{1} \unsure{of} $T$ which passes through $x',y'$} \deleted{and its transform} and \uline{compare} it \deleted{$\unsure{u(l_{y',x'})}$}. : $l^T_{u(y'),u(x')}$ (which is defined by the hypothesis of horizontality)
% The period belongs physically just after the cancelled lower formula;
% the separate colon is higher, after comparer.
\[
(53)\qquad \deleted{\unsure{u}}{}^u l^T_{y',x'}=g^T_{x,y}(u)\,l^T_{u(y'),u(x')}
\]
where
\[
g^{T\deleted{\unsure{u}}}_{x,y}(u)\in\widehat{\pi}_1\bigl(X',
\underset{\substack{\text{horizontal stratum}\\\text{of }X'_T\text{ which contains}\\u(x'),u(y')}}{
  \tikz[baseline=(term.base)]{
    \node[inner sep=0pt] (term) {$T(u x',u y')$};
    \draw[line width=.4pt]
      ([yshift=-1pt]term.south west) -- ([yshift=-4pt]term.south west)
      -- ([yshift=-4pt]term.south east) -- ([yshift=-1pt]term.south east);
  }
}\bigr)
\]
This being so, we shall have
\[
(54)\qquad \widehat f\underbrace{\bigl(g^T_{y',x'}(u)\bigr)}_{\substack{\rotatebox{-90}{$\in$}\\\widehat{\pi}_1(X,T)}}=g^T_{y,x}(u)
\]
\mspage{102}
